\section{Project diary}
\label{sec:diary}

The four-week development cycle was punctuated by eight merged pull requests and a
handful of direct-to-\texttt{main} commits for docs and trivia. The narrative below
follows the actual order of decisions, with commit hashes given so a reviewer can
verify any claim from \texttt{git log}.

\subsection*{Week~1 \textemdash{} Foundation (2026-04-25 to 2026-04-29)}

The first commit (\texttt{4ff85c1}) bootstrapped an empty repository. Four days
later, \texttt{ecdf3bb} (``Initial project scaffold for QuantBacktest'') landed the
full directory layout we would carry through the project: a \texttt{backend/}
package with \texttt{agents/}, \texttt{api/}, \texttt{analytics/},
\texttt{backtest/}, \texttt{data/}, \texttt{database/}, \texttt{llm/},
\texttt{strategies/} subpackages; a \texttt{frontend/} folder primed for Vite; a
\texttt{tests/} folder; and \texttt{requirements.txt} pinning the backend stack.
The agent boundaries were committed to at this stage rather than discovered later
\textemdash{} we used the breakdown in the rubric (data, strategy, backtest,
analytics, orchestrator, explanation) as the starting point and held to it.

\subsection*{Week~2 \textemdash{} First domain feature + first LLM call (2026-05-08)}

Two commits on the same day represent the project's two first proofs of life. PR~\#1
(\texttt{376d30e}, ``feat: add TSLA as new equity asset'') validated the end-to-end
data path: register an asset in the seed list, run \texttt{load\_initial\_data.py},
see daily bars appear in the database. \texttt{8997d20} (``feat: add LLM generated
report'') validated the LLM seam: we wanted to make sure the \texttt{LLMProvider}
abstraction could carry a real response before we built the rest of the analytics
layer around it.

We deliberately wrote the LLM feature \emph{first} as a worst-case probe. If it
worked, every other deterministic feature would be easier; if it failed, we had time
to redesign. It worked, and the rest of the project never had to compromise its
determinism for the sake of the AI report.

\subsection*{Week~3 \textemdash{} Frontend MVP, hourly bars, charts (2026-05-12 to 2026-05-16)}

Frontend Version~1 (\texttt{1f2f5c7}) landed on 2026-05-12. The same day, PR~\#2
(\texttt{1f13139}, ``Add native calendars and hourly bar support'') reshaped the
data layer: the OHLCV composite primary key became \texttt{(asset\_id, ts,
timeframe)}, so daily and hourly bars could coexist for the same asset, and the
\texttt{OHLCVCleaner} learned to reindex onto NYSE for equities, a bespoke
24$\times$5 window for FX, and 24$\times$7 for crypto. This was the project's most
intricate single change: a missed edge case in any of the three calendars would have
silently misstated the annualisation factor for every Sharpe ratio we report.

The chart suite then went through a flurry of small fixes (\texttt{3e22d4c},
\texttt{162ceba}, \texttt{ab1e2de}) that brought the strategy-category badges,
chart-tooltip alignment, and sidebar navigation up to standard. \texttt{a3b6635}
(``Chunk upserts, add LLM retries, and fix time indexes'') addressed the SQLite
parameter-count ceiling we hit when bulk-loading multi-year hourly history.

Late in the week, the PDF-export feature shipped (\texttt{0a7edb1}). The choice to
render the report client-side via \texttt{jsPDF} rather than server-side via a
headless browser kept the backend dependency footprint small and avoided the need
for the OS-level Chromium libraries that other PDF approaches require.

\subsection*{Week~4 \textemdash{} CI hardening, data layer, strategies, benchmarks (2026-05-19 to 2026-05-22)}

The week opened with two Claude Code GitHub Actions workflows: \texttt{claude.yml}
(\texttt{3d50b41}) for \texttt{@claude} mentions on issues and PRs, and
\texttt{claude-code-review.yml} (\texttt{6a38f60}) for an automatic
\texttt{/code-review} pass on every PR. The team's view of generative AI as a
collaborator-not-author shaped both workflows: they only respond to explicit human
triggers and they never self-merge.

Build-artefact hygiene then dominated two days (\texttt{4ce0311},
\texttt{820d610}). TypeScript compile outputs had begun leaking into
\texttt{frontend/src/}, which would have polluted every diff. PR~\#4 added the
\texttt{check-no-build-artifacts.yml} workflow that fails the build if any
\texttt{.js}, \texttt{.d.ts}, or stale \texttt{.tsbuildinfo} appears where it
should not.

The data layer consolidated on yfinance as the primary source across every asset
class (PR~\#5, \texttt{c2b6cce}). The Binance \texttt{ccxt} fallback remained for
crypto hourly windows older than the 730-day yfinance cap, and Stooq remained for
FX EOD fallback. Fixed-fractional sizing was corrected for crypto and FX (PR~\#6,
\texttt{ae082ef}).

PR~\#7 (\texttt{b0676ed}, ``added 6 new strategies'') doubled the strategy
catalogue to eleven by adding MACD Crossover, Ichimoku Cloud, Keltner Channels,
Time Series Momentum, Stochastic Oscillator, and CCI. Each was added to the
registry, exposed via JSON Schema so the frontend form rendered automatically, and
covered by a unit test that asserts the signal series on a known fixture.

PR~\#8 (\texttt{c0801c7}, ``Support benchmark overlays (API + UI)'') and its
sibling commits (\texttt{be6cf8a}, \texttt{55fa4b0}, \texttt{d432105}) brought the
chart suite to its final polish: a same-asset buy-and-hold and an SPY buy-and-hold
curve are now computed and persisted at run time, every chart shares a single
legend that toggles them, and a shared crosshair tooltip keeps values aligned
when the mouse hovers anywhere on the timeline.

The week closed with \texttt{d1efa3f} (Quick-Start docs) and \texttt{4a9e6ae}
(yfinance/httpx/curl\_cffi version bumps). The yfinance bump from 0.2.46 to 1.3.0
required adding \texttt{curl\_cffi} as a transitive dependency \textemdash{} a
small change that we noted in \texttt{CHANGELOG.md} because it could surprise
contributors who ran into install issues on locked-down corporate machines.

\subsection*{Documentation pass (2026-05-24)}

\texttt{79e200a} (``Updated all the current documentation files'') refreshed the
existing docs \textemdash{} \texttt{README.md}, \texttt{ARCHITECTURE.md},
\texttt{CLAUDE.md}, \texttt{ONBOARDING.md}, and the four
\texttt{docs/\{agents,calendars,data-sources,strategies\}.md} \textemdash{} against
the final code state. A second pass on the same day created the artefacts the
rubric requires (this PDF among them): \texttt{AGENTS.md}, \texttt{CHANGELOG.md},
\texttt{CONTRIBUTING.md}, \texttt{CITATIONS.md}, \texttt{AI\_USAGE.md},
\texttt{SECURITY.md}, \texttt{CODE\_OF\_CONDUCT.md},
\texttt{docs/user-guide.md}, \texttt{docs/api.md},
\texttt{docs/architecture-diagrams/}, and this \texttt{docs/academic/} bundle.

\subsection*{Post-submission CI automation (2026-05-24)}

On the same evening as the submission, PR~\#9 (\texttt{612e39e}, ``chore(ci): add
nightly headless docs-refresh workflow'') introduced \texttt{nightly-docs.yml}: a
GitHub Actions job that wakes at 02:00~UTC each night, runs Claude Code in headless
mode with a read-only Bash allowlist, and opens a pull request whenever it finds
documentation drift against the last 24~hours of commits. The workflow never touches
\texttt{main} directly.

Two small fixes followed immediately. \texttt{4dfc948} switched the authentication
method to \texttt{CLAUDE\_CODE\_OAUTH\_TOKEN} (the token type the
\texttt{claude-code-action} expects), and PR~\#10 (\texttt{58a4811}) granted the
\texttt{id-token: write} permission that the action requires to mint a GitHub OIDC
token at startup.

\subsection*{Node 24 Actions upgrade (2026-05-24)}

GitHub announced the retirement of Node 20--based runner actions on 2026-06-02.
Commit \texttt{01d9c30} (PR~\#12) bumped the three flagged actions
(\texttt{actions/checkout}, \texttt{actions/setup-python}, \texttt{actions/setup-node})
to their v6 majors across all four workflows in a single pass. The permissions comment
in \texttt{nightly-docs.yml} was updated to record why \texttt{id-token: write} is
required: \texttt{claude-code-action} mints a GitHub OIDC token at startup regardless
of which model-auth path is active, so the permission is not optional.

\subsection*{Things we tried and abandoned}

A few decisions did not survive the four weeks:

\begin{itemize}
  \item \textbf{Asynchronous backtest submission.} The first design queued runs and
        polled status via WebSocket. We backed it out because synchronous \texttt{POST}
        worked well enough at v1 scale and removed an entire layer of state to debug.
        The Dashboard still polls \texttt{GET /backtests} every five seconds so it
        will surface any future async-submission feature without UI changes.
  \item \textbf{Postgres in development from day one.} We initially set up Docker
        Compose with Postgres but it slowed local iteration noticeably. SQLite became
        the default and the schema was kept on SQLAlchemy generic types so a
        \texttt{DATABASE\_URL} swap reaches Postgres without code changes.
  \item \textbf{LLM tool-use via the provider SDK's native function calling.} The
        Orchestrator Agent now parses tool calls from a JSON envelope in the model
        output rather than relying on Gemini's native function-calling API, because
        the latter is not directly portable to other providers. The cost is a
        small parsing step; the benefit is provider-agnostic tool dispatch.
\end{itemize}
